% Figure floats for the two abort figures, written by figures/aborts.py.
% Every number in a caption is read off the replayed episode, so a changed
% tree cannot leave a caption asserting something the picture no longer shows.
\begin{figure}[tbp]
    \centering
    \thesisfig{abort_avoidable.pdf}
    \caption{The shape of a six-player enclosure abort. Qwen3.8, S3 board 10, \textsc{freeform}, move 57 of 59: B's only free neighbour is \texttt{R8C5}, and D plays into it although \emph{up} and \emph{left} were equally legal and would have left every player a move. Thirty of the thirty-two enclosure aborts recorded for the strong tier take this form; Table~\ref{tab:seal-greedy} gives what the sealing move bought the player that made it. Discs are objects and dark cells walls; a ringed disc has declared done and will not move again.}
    \label{fig:abort-avoidable}
\end{figure}

\begin{figure}[p]
    \centering
    \thesisfig{abort_forced.pdf}
    \caption{The other two, and the line that avoids them. The same board as Figure~\ref{fig:abort-avoidable}, sealed identically under \textsc{none} and \textsc{structured}. The top row is the recorded play: D's move at 38 leaves E a single legal move, and at 39 that move lands on the one cell B had left. The bottom row replays the same two plies with the other legal choice at each, after which every player still has a move and B keeps its 4-move route home, drawn dotted. Each panel shows the board \emph{before} the move drawn on it.}
    \label{fig:abort-forced}
\end{figure}
